%-*-tex-*-
\ifundefined{writestatus} \input status \relax \fi %
\chcode{pagsiz}

\def\cqu{\cquote{It is when I am struggling to be brief that
I become unintelligible.}{Ars Poetica, Horace (65-8~BC)}
}

\chapterhead{pagsiz}{PAGES:\cr SIZES\cr MAGNIFICATION}
\intex\ has a number of page size parameters that are set to make the output
possible on different devices, and to facilitate special formats. 
All page size
parameter changes are made in assignment fashion with an optional |[=]| sign.
All page size changes must be made inside a 
\@|\documentstyle| list in order to insure that they take effect.
Figure~\ref{pagdiag} diagrams the various dimensions.

\input pagdiag 
\beginpageinsert
\figureform{pagdiag}{Diagram of the various page dimensions.\captionbody{
     The actual printing area is determined by the computed dimensions
{\tt \1hinnerpagesize} and {\tt \1vinnerpagesize}. The {\tt outer} dimensions
are normally specified. The {\tt \1firstcolumnoffset} is normally
0pt.}}{\centergraph{ \pagdiag sc:.8 }}
\endpageinsert

\shead{pagecomlist}{Command List}
\beginthreecolumn
\hfuzz=20pt
\ext\@|\documentstyle|
\ext\@|\houterpagesize|
\ext\@|\hpapersize|
\ext\@|\leftmarginsize|
\ext\@|\leftmarginskip|
\pri|\mag|
\pla|\magstep|
\ext\@|\topmarginsize|
\ext\@|\topmarginskip|
\ext\@|\vouterpagesize|
\ext\@|\vpapersize|
\endthreecolumn

\shead{pagesandmargins}{Page Size Setting and Margins}
These commands define and set the various page sizes. Only the |outer| and
|header|, and |paper| commands should be changed. 
\intex\ also has some |inner| and |column| sizes. These are computed and
should not be modified -- although they are accessible. See
Section~{\ref{inner}}, page \pref{inner} for details. 
{\bf All dimension changes must be made inside a |\documentstyle| command or
they  will not necessarily take effect.}

This  section details the various parameters used to
name and describe the parts of the printed page that can be varied.
The general convention is that |\h<name>| names a horizontal dimension
and |\v<name>| names the corresponding vertical dimension. The various 
|<name>|s are
\bshortcomlist
|outerpagesize|&{this is the extent of the area allowed
                                    for printing. It includes all headers,
                                    footers, footnotes \dots\
                                    \@|\vouterpagesize| and \@|\houterpagesize|
                                    are the vertical and horizontal 
                                    forms}\cr
|papersize|&{this is the ``physical'' size of the paper and includes all 
                                   the |outerpagesizes| plus 
                                   margins. The actual 
                                   vertical extent of the paper 
                                   is  \@|\vpapersize| and
                                   the actual width of the 
                                    paper is is the size of \@|\hpapersize|.
The default values are correct for 
11in by 8.5in paper but must be increased for 
14in by 8.5in (Legal size) paper.}\cr
\@|\vsize|&the actual verticalsize of the present page. This is set at the time when the
first material is submitted for putting on the page. In \intex\ it is usually
the |\vcolumnsize|. It should be changed only with knowledge.\cr
\@|\hsize|&the actual horizontal size of the present {\bf box} in which \tex\ is
working. This may be a page or some box that is to be eventually placed on
the page. 
This is set at the time when 
first material is submitted for putting on the page. In \intex\ |\hsize|  is
usually 
the |\hcolumnsize|. \cr

\eshortcomlist

Assignments are made in the normal way. The general form is

\@|\<dimen name> [=] <dimen>|
\noindent
where |<dimen>| is a dimension. Thus the papersize, margins, 
and outer page sizes are set in \intex\ by 
\begintt
\documentstyle{\vpapersize = 11in
               \hpapersize = 8.5in
               \vouterpagesize=9in 
               \houterpagesize=6.5in
               \leftmarginsize=1.25in
               \topmarginsize=1in}
\endtt
These values are actually set for the paper style in {\tt papersty.tex}. 

\shead{magnification}{Magnification}
\tex\ is able to magnify an entire document through its
|\magnification| parameter. The form is
\beginblockmode
\mbr
\pri\@|\mag [=] <1000 times scale factor>|
\nbr
The document is magnified by the |<scale factor>|, times any  |texprint|
magnification. It can appear only once and should appear at the beginning of
the document. {\bf It should be used with extreme caution. The results may 
be too big for the output device.} The {\it plain} form |\magnification| has
been defined to be the same as |\mag|. The reason is that |\magnification| in
{\it plain} also sets page sizes which is inappropriate here. \TeX Graph will
scale correctly when using |\mag|. 

\mbr
\pri\@|\magstep<integer>|
\nbr
This is a valid magnification scaling for the |\magnification| command. The
actual magnification is $1.2^{\langle{\rm integer}\rangle}$. 
Thus |<integer>=3| implies a
magnification of 1.728.
\endblockmode

In fact there are two points at which
magnification can occur. The first is in the \tex\ document and the second is
when it is printed. All that is important is the final magnification. A
magnified document may be printed only if all the fonts are available at the
magnifications requested. In general, magnification should be controlled by 
\intex\ so that \TeX Graph will operate correctly. 

\shead{inner}{Internal and Computed Sizes}
\intex\ is actually in multicolumn mode at all times. The default though is
just one column. In order to live in this world \intex\ understands |inner|
and |column| sizes. The |column| sizes are the actual sizes that are used to
for the text. The |inner| horizontal size is the sum of the column widths and
the intercolumn spacing. To this is added a \@|\leftcolumnoffset| (twice) to
obtain the |\houterpagesize|. Normally the |\leftcolumnoffset| is zero. The
effect of a left column offset is to inset the pages with respect to the
header/footer boxes whose width is the |\houterpagesize|. Thus 

\bshortcomlist
|columnsize|&this is the actual size of the column
                                    where the text is printed. Footnotes 
                                    are inside the column. In single 
                                    column format columnsize and innerpagesize
                                    are the same 
                                    (\@|\vcolumnsize|, \@|\hcolumnsize|) \cr
|innerpagesize|&this is the size of the page left after
                 the headers and footers have been removed 
                 (\@|\vinnerpagesize|, \@|\hinnerpagesize|) \cr
\eshortcomlist
\ejectpage
\done
